\begin{table*}[htbp]
\centering
\small
\setlength{\tabcolsep}{4.5pt}
\caption{Comparación cuantitativa de los resultados de simulación computacional: coeficientes de amortiguamiento magnético $b$, velocidad terminal asintótica $v_t$, tiempo total de tránsito $t_{\mathrm{caída}}$ a lo largo de $1.00\ \mathrm{m}$, y balance de incertidumbres por análisis analítico y método de Monte Carlo ($N=1000$).}
\label{tab:resultados}
\begin{ruledtabular}
\begin{tabular}{lcccccc}
\textbf{Material del Tubo} & $b\ [\mathrm{N\cdot s/m}]$ & $v_t\ \text{(Analítico)}\ [\mathrm{m/s}]$ & $v_t\ \text{(Monte Carlo)}\ [\mathrm{m/s}]$ & $t_{\mathrm{caída}}\ [\mathrm{s}]$ & $E_{\mathrm{Joule}}\ [\mathrm{mJ}]$ & Error $\delta v_t / v_t$ \\
\hline
Cobre Electrolítico & $3.053 \pm 0.251$ & $0.0915 \pm 0.0076$ & $0.0924 \pm 0.0080$ & $10.97 \pm 0.94$ & $282.2$ & $8.3\%$ \\
Aluminio 6061-T6 & $1.931 \pm 0.160$ & $0.1447 \pm 0.0120$ & $0.1461 \pm 0.0127$ & $6.98 \pm 0.60$ & $281.2$ & $8.3\%$ \\
PVC (Caída Libre) & $\approx 0$ & $4.4287 \pm 0.0044$ & $4.4287 \pm 0.0044$ & $0.463 \pm 0.003$ & $\approx 0$ & $0.1\%$ \\
\end{tabular}
\end{ruledtabular}
\end{table*}
